When we study physics beyond doing long calculations, formulate complicated theories or doing extremely elaborated experiments in an advanced laboratory there is a really important fact many people ignore, the \textit{story of physics} if we want to study physics we must have understand where was born the physics and more important, how was it done. Here we are going to check the story of physics from the geometrical optics, to corpuscular optics to wave optics, to electromagnetic optics and the most recent the quantum optics.

    \section{Optics from Euclid to Huygens}
    
         Here I'm going to copy and paste one of the best articles I've read about the story of optics \citep{herzberger1966optics}.\\

         If I were endowed with dictatorial powers, I would require everyone receiving a degree in a scientific subject to know its history and to have read the classical papers relating to it.
         Historical knowledge is important because it stimulates creative thinking. The man who first struggled with an idea, trying to find a law, looked at the situation with different eyes than do we who accept the law as a matter of course. He considered alternatives to the law, and different interpretations, and some of these alternatives and interpretations may still be stimulating and worth thinking about.\\
         
         I was fortunate during the eight years of my affiliation with the Zeiss works at the beginning of my career to have access to their extensive library; it contained the originals, or at least translations, of nearly all the writings in optics from Euclid to contemporary times. Over thirty years ago, I wrote a historical survey describing the content of important papers on geometrical optics. This task was immensely facilitated by the extensive bibliography in the third edition of Czapski-Eppenstein,I to which I had contributed my small part as a work student at the firm of Zeiss. Lately, I have looked through most of these papers again, as far as they were available to me, and have discovered many important details that I had not understood at that time. I should like to share these with you. I emphasize that my interest is not basically historical. I am not interested so much in questions of priority and originality of documents. My object is to entice you to read and to enjoy the old scriptures, and maybe take along an idea or two that might be of interest.

         \subsection{Ancient times}

            I am sure that the Egyptians and, even more, the Babylonians, with their great interest in astronomy,must have had some knowledge of optical laws. But the first important mention of one of the phenomena of refraction I found was in Plato's Republic.' Plato(427-347 B.C.) quotes Socrates as saying (Jowett's translation): `` . . . the same object appears straight when looked at out of the water, and crooked when in the water; and the concave becomes convex, owing to the illusion about colours to which the sight is liable. Thus every sort of confusion is revealed within us; and this is that weakness of the human mind on which the art of conjuring and of deceiving by light and shadow and other ingenious devices imposes, having an effect upon us like magic.¨ Our senses give an erroneous picture of reality; they are untrustworthy;alone can lead us to truth pure thinking. This idea persisted for two millennia under the influence of the misinterpreted teachings of Aristotle.\\

            How we see is a question that long occupied the Greek philosophers, and, through them, the Neo-Platonists and the Scholastics. (Indeed, it cannot be said that the question is settled, and in recent times Ronchi  has set forth some cogent ideas.) Even before Plato there was speculation in the writings of the Greek philosophers about the process of seeing. Most of the Greek thinkers believed that the eye sends out rays to the object, and Aristotle (384-322 B.C.) seems to have been the first to question this belief. If, says Aristotle, the eye sends out fiery rays that stream from it as from a lantern, why can we not see in the dark? His explanation, that the eye contains a transparent substance and that seeing is associated with the motion of something between the object and the eye, seems quite modern.\\

            After the conquest of Greece by Alexander, Alexandria (in Egypt) and Sicily became the strongholds of Greek culture. The first great textbook of optics was written by Euclid of Alexandria, who lived about 300 B.C. And is better known for his Elements of Geometry. This famous work is still must reading for every scientist and is available in several editions. When King Archelaos once complained about the difficulty of learning mathematics and asked whether there was not an easier way, Euclid's answer is supposed to have been: ``There is no royal route to geometry¨. When one of his pupils asked him what of material value could be gained by learning geometry, he told a slave, ``Give the man threepence and tell him to go away¨.\\

            Under Euclid's name we have two books of interest in the present connection: Optics and Catoptrics. The first deals with the laws of vision and perspective and the second with reflection. An English translation of the first was published in the Journal of the Optical Society of America twenty years ago. 7 It is a striking fact that one of the definitions in this book contains an idea that could be used to develop many problems of diffraction theory.\\

            I quote some of Euclid's definitions, slightly paraphrased from the Optical Society of America translation: ``Let it be assumed that lines drawn directly from the eye pass through a space of great extent [not infinite]; that the form of this space is the mantle of a cone with its apex in the eye and its base at the limits of our vision; that those things upon which the cone falls are seen and those things upon which the cone does not fall are not seen; and that things seen through several cones simultaneously appear to be more distinct.¨ From this, Euclid concludes, among other things, that nothing can be seen in its entirety. The same objects located nearby are seen more distinctly than if they were at a distance. Every object that is visible has a certain limit of distance. If that is reached-that is, if the object lies entirely inside the mantle of the cone of vision-it is no longer seen. These are the axioms from which he derives the laws of perspective.\\

            In Catoptrics the law of reflection is stated, namely, that incoming and outgoing rays form the same angle with the surface normal. From this he derives the exact position of the image behind a plane mirror and the principle that left- and right-hand sides are interchanged. He generalizes these laws for spherical mirrors, both concave and convex, describing virtual images as well as stating that concave mirrors can give real images. As the position of the image is formed by reflection at curved surfaces, he suggests the point of intersection of the ray with the axis (shades of my definition of diapoint!). He shows the existence of a focal point for spherical mirrors, although he errs with respect to its position. One of his theorems says that an eye at the center of a spherical mirror sees only itself. Another states that you can light a fire at the focus of a concave mirror. He states that straight lines appear convex in a convex mirror. One can see the same object through many plane mirrors if they are positioned along the sides of a regular polygon with two more sides than there are mirrors.\\

            According to Roman historians, the greatest mathematician of ancient times, Archimedes (287-212 B.C.), is supposed to have used mirrors to destroy the Roman fleet besieging Syracuse. Like most scientists, I have doubted this story, but a personal experience made me doubt my doubts.\\

            Two years ago I had an optical accident. My wife had me get a large bottle of distilled water, and I left it in my auto in front of our laboratory. Driving home that evening I smelled smoke and stopped at a garage. After thoroughly investigating all the mechanical details and finding them in order, the mechanics discovered that smoke was coming out of the upholstery. The cylindrical bottle had acted like a lens and the sun had made a hole in the upholstery and ignited the stuffing inside. After this personal experience, I am somewhat more reluctant to doubt the story of Archimedes' mirrors.\\

            The next person I should like to discuss, Hero of Alexandria, lived sometime between about 150 B.C. and A.D. 250. We have two outside dates for this philosopher: he mentions Hipparchos, who flourished between 146 B.C. and 126 B.C., but the mathematician Pappus (ca. A.D. 300) follows Hero's treatment in his own writings on mechanics. Be this as it may, Hero's Catoptrics is one of the most interesting books on optics. He is not content with Euclid's statement that light travels in straight lines, nor is he content with the form of the reflection law. He looks for a reason for both laws and derives it from the assumption that light chooses the shortest possible path between two points. We know that this minimum principle is not quite correct: the path is, in general, only an extremum. It may or may not be significant that Hero's book illustrates the concept with a picture of a convex mirror, for which the minimum principle is correct-not a concave mirror, for which it can be wrong. We shall find the same trick worked later by Fermat, who extended Hero's law to derive the refraction law. However, Fermat was called on the carpet for this by Abbé Mersenne, one of his contemporaries.\\

            Hero of Alexandria-who, by the way, in his Mechanics described the so-called seven simple machines known to us-not only discussed spherical mirrors in his Catoptrics but, as an early practical experimentalist, also showed what can be done with cylindrical mirrors. If you go into an amusement parlor and see yourself in mirrors of all kinds (I especially like those that elongate the height and diminish the girth), you know the proprietor has read Hero of Alexandria. If you see a play in a tanagra theater, where the actors appear on the stage, by the skillful use of mirrors, as beautifully dressed dwarfs, you know that Hero of Alexandria wrote the prescription. If you want to see around the corner or look at a person without letting him see you, Hero of Alexandria tells you how to do it with mirrors.\\

            Ptolemy of Alexandria, who lived about A.D. 150, is best known for his Almagest, but he was also the first to make a serious attempt to study the law of refraction. He wrote a five-part volume entitled Optics, several copies of which have come down to us. This work, of which the first part has been lost (although we have a description of its contents), can be summarized as follows.\\

            The first part deals with the relations between the eye and light, and the second with the conditions of visibility: size, form, color, and motion of bodies. Ptolemy explains why we usually see only one image of an object, although we have two eyes, and describes the conditions under which we see more than one image. The third part deals with the laws of reflection by plane and convex mirrors. Ptolemy gives probably the first attempt at an explanation of why the moon and the sun appear larger at the horizon than at the zenith a question that is still being discussed.

            In the fifth part, Ptolemy tries to find a law for refraction. First he observes that for small incident angles the angles of incidence and of refraction are proportional to each other. He also remarks that, in transition from a thinner to a denser medium, the refracted ray comes earer to the perpendicular, the opposite being true for the reverse direction of travel. He uses an ingenious contraption to measure the incident and the refracted angles for the transition from air into water, air into glass, and water into glass. This instrument consists of a circular disk divided into 3600 with a colored marker in the center and two movable markers on the periphery. You put the instrument into water, line up the three markers, and read the two positions, one giving the incident angle and the other the refracted angle. Ptolemy gives refraction tables in 100 intervals. My friend, the late Hans Boegehold, studied hese tables carefully and found that the values for 500 and 600 incident angles are surprisingly accurate, but that the other angles were obviously calculated by a second-degree interpolation formula. (The first differences are exactly 1 30'.14) Ptolemy applies the same method for the transition of light from air into glass and from water into glass. He also discusses atmospheric refraction, the cause of which he assumes to be that the densities of air and of ether vary with height. He points out that only for a star in the zenith do the true and the apparent positions coincide.

            
        \subsection{The arabain period}

            After the fall of Greece, scientific progress ceased on the continent of Europe for nearly a century and a half, but Alexandria remained the center of scientific discovery. When the Arabs took over, they brought their scientific inheritance and knowledge of mathematics and the mathematical sciences into northern Africa and Spain. It is remarkable that, at a time when magic and mysticism and misinterpretation of Aristotelian doctrines prevented scientific progress over most of Europe, the Arabs held high the torch of science in their countries.\\

            The next big contribution to our science was by Ibn al-Haitham, or Alhazen of Cairo, whose date is frequently set around 1100, although some sources tell us he lived from 965 to 1038. Alhazen's Optics" consists of seven parts. The first gives an anatomical description of the eye, which, he maintains, contains three liquids separated by four membranes. The names which he gives to them the watery medium, the crystalline medium, and the glassy medium are still used today. The phenomenon of binocular vision is explained by assuming that the pictures from the two eyes are transported to one visual nerve. Alhazen discusses the pupil of the eye and speaks of a pyramid of rays (rather than a cone) coming from the object to the eye.\\

            The second part, like the corresponding book by Ptolemy, deals with the properties of bodies, and he distinguishes no fewer than twenty-two different properties.
            
            The third part deals with optical illusions. He draws attention to the illusions arising from a false interpretation of a picture by the human mind and imagination. The fourth, fifth, and sixth parts deal with problems of reflection and mirrors. One of the problems studied in great detail has since been called after Alhazen: it consists in finding the point on a mirror where the incident and the reflected rays meet. Alhazen corrects errors of the Greeks about the position of the images in spherical, conical, and cylindrical mirrors. He also emphasizes the fact that the incident and reflected rays lie in the same plane.\\

            Alhazen's seventh part deals with refraction. He proposes methods for measuring incident and refracted rays through different media, but he gives no tables. He knows that a glass sphere magnifies. He explains the apparently large size of the sun at the horizon by the assumption that we compare the sun with objects of known size lying along the line of sight. He investigates binocular vision and makes a good guess at the height of the atmosphere. He breaks with the idea of the Greeks that light rays originate at the eye and go to the objects, supporting his claim by the circumstance that sunlight can hurt the eye and that we can see afterimages.\\

            Other manuscripts show that Alhazen used the pinhole to view solar eclipses and that he knew how the mage becomes blurred and finally becomes circular as the pinhole increases in size. He did not write as though this were a discovery of his own, and Aristotle  had noted that the spot formed by the rays from the uneclipsed sun is circular when the rays pass through a rectangular hole in wickerwork if the hole is small, but becomes rectangular when the hole is large. However, Alhazen's is the first known unequivocal description of the camera obscura. Alhazen's works were of enormous influence in the development of optics.

        \subsection{The middle ages}

            The knowledge of the Arabs was transmitted to Europe via Spain and Sicily and set in motion the intellectual revival that culminated in the Renaissance. The first of the European scientists to play a role in optics was Vitello, a Polish mathematician who lived and studied in Italy about 1270. His book in ten parts is probably the most voluminous optical work ever written. It contains about five hundred closely printed folio pages. Unfortunately, he is more famous for his errors than for his contributions.\\

            He reissued the tables of Ptolemy and added tables for the passage of light from water into air, from glass into air, and from glass into water. It is unfortunate that he never tells us how he arrived at his values, for they include impossible data, being in contradiction with the then unknown phenomenon of total reflection.
            
            It was this error that later prevented Kepler from finding the correct form of the refraction law. Nevertheless, Vitello made two notable contributions. He explained the colors of the rainbow as originating from refraction and reflection by the water drops, and he made the proposal to collect the light of the sun into a point by using parabolic mirrors.\\

            The Englishman Roger Bacon (1215-1294) discovered the spherical aberration of a mirror; he further stated that rays forming a cone around the axis in object space also form a cone around the axis in image space. From the magnifying property of a sphere, he speculated that it might be possible to invent instruments in which a man would look like a mountain and a small group like a big army. However, since he did not say how this could be achieved, it is doubtful whether he should be called (as he was by Sir John Herschel) the inventor of the telescope. His writings are so sketchy and in definite that it is impossible to attribute specific inventions to him or to reach any conclusions about what he actually knew.

        \subsection{The renaissance}

            Coming to the Renaissance, we must mention Leonardo da Vinci (1452-1519). His notebooks show that he studied the eye anatomically and made a model of it, although he supposed the rays to cross twice in it to form an upright image. 18 , 9 Immediately afterward he described the camera obscura but correctly showed the image inverted!20 He described, with sketches, a Rumford-type photometer and a machine for grinding concave mirrors. Unfortunately for his optical reputation and the development of optics, he never organized his notes into publishable form, so they were practically unknown until Venturi published them in 1797.\\

            About half a century later came another Italian we should mention, Francesco Maurolico (1494-1577), who solved one of the problems that had bothered previous thinkers: that the image of the sun is round even if formed by a square pinhole. He improved our knowledge of vision by suggesting that the crystalline lens of the eye works like a double-convex lens. This enabled him to attribute myopia and hyperopia to an eye whose lens has a larger or smaller curvature than is proper for the length of the eyeball. He emphasized for the first time that the rays from an object point envelop a surface that is now called the caustic surface and that this surface is the place of the highest light concentration. And he explained the colors of the rainbow by multiple reflection in the water drops. All this information was concisely organized by 1567 into a work entitled Photismi de lumine , although it was not published for some thirty years (1611) after the author's death, and then only as a consequence of the interest in optics engendered by Galileo's telescopic discoveries.\\

            In Italy, under the influence of the philosophy of the time, scientific progress and magic were mixed in a way that led to interesting reading. The development of mathematical methods, however, served to gradually separate speculation and imagination from the solid basis of observation. 
            
            One of the most colorful books of the period was the Magia naturalis of Giovanni Battista della Porta (ca. 1538 or 1543 -1615). In his youth, della Porta had written a four-part volume entitled De miraculis rerum naturalium (1558) and, subsequently, many books and plays, including De refractione (Fig. 6) in nine parts in 1593. The Magia naturalis, which was an expansion of the earlier De iraculis into twenty volumes, appeared in 1589. It became so famous that it was soon translated from Latin into the vernacular languages Italian, French, Spanish, and English. It is a strange medley. Only the seventeenth part concerns optics, but some of the others are also worth noting for our amusement if nothing else. I mention the titles of a few of the twenty parts, to give you an idea of their contents. The first is entitled, "The Causes of the Miraculous. It deals with astrology, sympathy and antipathy, and the influence of place and time on events. The second deals with the origin of animals. He quotes many authorities who agree that snakes are made from the marrow of men and the hair of women. The third deals with garden plants; the fourth, with how to become wealthy by good housekeeping; the fifth, with how to make precious stones artificially; and the ninth, with make-up and other ways of improving the beauty of women (as if this were necessary!).\\

            Finally, the seventeenth part deals with focal mirrors and other mirrors and explains what one can do with them. This book exhibits progress in the understanding of optical phenomena. In it are discussed the same tricks that Hero had dealt with 1400 years earlier, namely, how to use mirrors to obtain desired effects. However, we find here the first exact theory of multiple mirrors, the first complete description of the camera obscura, with pinhole or lens, and a discussion of combinations of positive and negative lenses designed to improve vision. Porta compares the eye and its pupil to the camera obscura and its stop. (Leonardo had anticipated him here, but less definitely and accurately.) At the end of this century also falls the invention of the telescope and the microscope. \\
            
            Eyeglasses had apparently been invented toward the end of the thirteenth century, but they seem not to have attracted the attention of writers on optics until the practical Leonardo discussed methods of making them. Some lens grinder-more likely, several independentlymust have discovered the principle of the Galilean telescope by accident, and certain passages in Leonardo's notebooks a hundred years before Galileo can be interpreted as instructions for making such a telescope. Be that as it may, a Dutchman, Zacharias Jansen (1588-1632), was said by his son to have made a telescope in 1604 according to the model of an Italian dated 1590, and soon telescopes became well known in Holland. Galileo Galilei (1564-1642) heard of them and, apparently independently, made an instrument that he exhibited in 1609, and in his lifetime made over a hundred. The Dutch instruments were so imperfect that a magnifying power of more than 3X was impractical. Galileo made instruments with a useful power of approximately 30 X, an order of magnitude greater. He was by no means the inventor of the telescope or even the first scientist to describe it; Vavilov reminds us that Leonardo may have anticipated him and that Porta described, three months after Galileo's demonstration, an instrument he had made independently. But Galileo is outstanding on two counts: he revolutionized man's ideas of the uni- verse by his applications of the telescope and, by improving it by an order of magnitude, Ronchi points out that Galileo is entitled to be considered the father of the art of fine optics.\\

            Moreover, from telescope to microscope is a short optical step, and by 1624 Galileo was using lenses of short focal length and thus making precursors of the modern compound microscope. 9 A few decades thereafter, the Dutchman Antonj van Leeuwenhoek (1632-1723) improved the quality of simple lenses used for magnifiers or simple microscopes, and in a letter of 1674 he began the descriptions of the discoveries for which he became justly famous.24 Meanwhile, Robert Hooke (1635-1703) in England also improved the quality of lenses in order to make the compound microscope practical, and in 1665 the Royal Society published his Micrographia.\\

            Aside from publications, which were in an embryonic state, the seventeenth century brought a lively exchange of ideas and thoughts between scientists, mostly friendly but sometimes polemical and frequently expressed in code to conceal them from the uninitiated. It is therefore necessary to study not only the books of these scientists but also their correspondence (cf. Leeuwenhoek's letters quoted by Dobell ).
            
            The same period witnessed the development of scientific knowledge to the system that is still important.Indeed, scientific development in the seventeenth century was greater than in any century before or since. Scientific giants like Kepler, Descartes, Pascal, Fermat, Newton, Leibniz, and Huygens appeared and left their marks in all fields of science, including optics. Let us sketch some of their contributions to geometrical optics. \\
            
            Johannes Kepler (1571-1630) wrote two books on optics: Supplements to Vitelo in 1604 and Dioptrice  in 1611. The object of the first is to find the law of refraction. Kepler nearly succeeds, though he does not look for a relation between the incident angle and the refracting angle but between the incident angle and the deviation. He asks himself what form a surface should have to refract a parallel beam of light exactly into a point image. He studies the hyperboloid, which, as we know, fulfills the condition; but he rejects it because it disagrees with the values that Vitello gave for the relation between the angles of the incident ray and the refracted ray. If only the Kepler of 1611 and the Dioptrice, in which he discovered the law of total reflection, had reread his book of 1604, he would have seen that the data of Vitello had been erroneous and thus he might have found the correct form of the law of refraction.\\

            Kepler investigated the function of the eye more profoundly than his predecessors. He discovered that the eye is not at rest but moves around a point inside it, and he investigated the importance of this motion on our vision. He drew attention to the fact that the image on the retina is inverted. All this is in his Supplements to Vitello.\\
            
            His Dioptrice , although a relatively small book, was of importance to the development of optics. In it, Kepler assumes that for small angles the angle of refraction is proportional to the angle of incidence. From this assumption, he develops the laws of firstorder optics for thin lenses and systems of thin lenses with finite separations. He describes, among other things, the Keplerian telescope, in which two positive lenses are placed so that the back focus of the first is at the front focus of the second, and compares it with the form named after Galileo, consisting of a positive and a negative lens placed so that the focus of the former s at the virtual focus of the latter. He gives the focal length correctly for all forms of lenses. His book contains many practical proposals for lens design and is highly recommended to those who are interested in these problems.\\

            Rene Descartes (1569-1650)  must be considered as one of the discoverers of the correct law of refraction, for he published the sine law in his Dioptrique in 1637. After his death, he was violently attacked by Isaac Voss, who accused him of stealing the law from Willebrord Snel* (1591-1626). Descartes had indeed visited Leyden in 1630, where Snell, a professor of physics, had written a book on optics containing the refraction law. Although the book had never appeared, Voss, in his posthumous attack, claimed that Descartes had seen it on his visit after Snell's death, and Voss quoted Snell's proof.\\

            The question of priority has been much discussed in the literature. A recent study has been published by Boegehold, who refers to a paper by P. Kramer. Boegehold found that Descartes, in some letters written in 1627, had suggested making plano-convex lenses with hyperbolic surfaces. Since such a lens would give a sharp point image for an infinite object if and only if one assumes the sine law, it is fair to conclude that Descartes knew of the law three years before his visit to Leyden in 1630. It is strange that Voss should have years after waited for twenty-five years-twelve Descartes's death-to make his accusation. Credit for discovering the law of refraction should therefore be divided between Descartes and Snel, with a fair share going to Johannes Kepler, who failed to find the sine law only because he trusted the tables of Vitello.\\
             
            Descartes's Dioptrique is beautifully written. It tries to explain the law of refraction by means of the theory of elasticity. A light ray is compared to the path of a ball, first being bounced off the ground, then being slowed down as it goes into a denser medium. Descartes was the first to suggest determining the velocity of light from the position of the moon at an eclipse.\\

            He also made experiments with human and animal eyes. If the eye is removed and put in the opening of a window in a darkened room, one can see the images of distant objects on the retina, and by pressing the eye a little, one can see the image of near-by objects, showing that the eye can accommodate.\\

            In the eighth chapter of his book, Descartes investigates surfaces that image a point sharply-the socalled Cartesian surfaces, of which the conic sections constitute a special case. He proposes machines for grinding lenses and demonstrates the colors of the rainbow with small glass spheres filled with water.\\

            But much of Descartes's reasoning in connection with the refraction law was faulty and was rejected by the great mathematician Pierre de Fermat (1601-1665), who pointed out that water would have to be less dense than air to fit Descartes's explanation. Fermat returned to the formulation of the refraction law given by Hero, deriving it from the assumption that light has different velocities in different media and takes the quickest path between any two points. Like Hero, he illustrates this idea with a drawing of a convex surface, for which the minimum principle is correct, rather than a concave surface, for which it can be wrong. But he was called on the carpet by Abbé Mersenne, as was mentioned earlier, one of the great scientists and letter writers of his day. Fermat's answer is a classic example of double talk: he says that the path would be a minimum if the surface were a plane, and the incoming ray sees so little of the surface that it cannot distinguish whether the surface is plane or curved!

        \subsection{Newton and Huygens}

            Sir Isaac Newton (1642-1727), in his Lectiones opticae of 1669-1671 (published posthumously), showed how to calculate the spherical aberration of a plano-convex lens for an infinite object point. This book contains the formula for locating the image, still known by his name, in terms of the distances of the object and the image from the focal points. It also contains a thorough investigation of reflection in prisms, including the position of minimum deviation. It even goes into some aberrations, in particular, image errors as functions of color, and it compares the size of the color aberration with the spherical aberration of a given lens and prism.\\

            Although Newton's best known work in optics is his Opticcs , I should like to draw your attention especially to his beautiful paper on color, in which he describes the separation of white light into the colors of the spectrum by means of prisms. This charming paper should be read by every student of optics, especially since the writer not only gives his final conclusions but describes many false attempts at interpreting his experiments and the refutation of these attempts.\\

            Newton was misled into believing that all materials have the same relative dispersion, since he found this to be true for such different media as glass and water. The conclusion that he drew-that telescopes should not contain refracting surfaces-was erroneous. But it led him to study mirror telescopes and, with his teacher Gregory, he contributed much to the development of these instruments.\\

            Christian Huygens (1629-1695) wrote two books on optics: Dioptrical in 1653, published posthumously , and Traite de la lumiere , written in 1678 but not published for twelve years. The first book shows enormous advances with respect to the understanding of optical image formation. The second proposes a new and revolutionary theory of light, deriving refraction and reflection laws from a completely new principle.
            
            This principle, set forth earlier in the Dioptrica, is a starting point from which all the laws of Gaussian optics can be derived for lens systems with finite thicknesses. It is expressed as follows: if we interchange object and eye in an optical system, the object will have the same apparent size and orientation as before. The law is known to us in the form that the product of lateral magnification and angular magnification is constant in an optical system, and equal to the ratio of the refractive indices.\\
            
            Huygens found, before Newton, the law that, for a single surface or for a single lens, the image distance is proportional to the object distance when the distances are measured from the respective principal foci. He knew of the aplanatic surface. He was acquainted with the important property of the optical center of a lens (the point which divides the distance between the vertices in the ratio of the radii), namely, that a ray through it emerges parallel to its original direction. He had an approximate formula for the spherical aberration of a thick lens that enables the form for minimum spherical aberration to be determined.

            In his Traite de la lumiere, which should be required reading for every physicist and is now readily available, Huygens derives optical laws from the fact that the light emerging from an object point P forms a wave surface S at any given time t. The light rays are then the normals to this surface and are not straight unless the velocity is constant throughout the medium.\\

            The surface can also be constructed in the following way. At an arbitrary time , light from the point source P forms a wave surface S. (assumed to be spherical in the figure). Considering each point of S, as a source sending out a new wave, we see that the envelope of all the waves at the time gives the original wave surface S.
            
            With the help of this concept, Huygens derives the straight paths of light in a medium of constant velocity. He derives not only the reflection and refraction laws but also the law of refraction in the atmosphere and in doubly refracting crystals. He clearly defines the caustic as the evolute of the rays and investigates it after reflection and refraction, in addition to investigating the degenerate form of the wave surfaces in the points of the caustic. He proves Malus' law: that rays form a normal system for a single reflection and refraction.\\

            It should be emphasized, however, that this formulation is mathematically identical with ray optics, as Hamilton showed in his paper of 1831. The assumption that light has a finite wavelength and that it can be described as a sine wave going along a ray is foreign to Huygens. However, once this theory was established, the methods of Huygens enabled Fresnel to derive all the laws of diffraction in a most elegant manner.